\documentclass[12pt,a4paper]{article}
\usepackage{ctex}
\usepackage{amsmath,amssymb}
\usepackage{graphicx}
\usepackage{float}
\usepackage{listings}
\usepackage{xcolor}
\usepackage{geometry}
\geometry{left=2.5cm,right=2.5cm,top=2.5cm,bottom=2.5cm}

% 代码高亮设置
\lstset{
    language=Verilog,
    basicstyle=\ttfamily\small,
    keywordstyle=\color{blue},
    commentstyle=\color{gray},
    numbers=left,
    numberstyle=\tiny,
    frame=single,
    breaklines=true
}

\begin{document}

\title{\textbf{定点乘法器性能优化实验报告}}
\date{\today}
\maketitle

\section{实验目标}
\begin{enumerate}
    \item 理解原始32位串行乘法器的工作原理（32个时钟周期完成运算）。
    \item 对乘法器进行最简单有效的性能优化，将时钟周期从32周期减少至16周期。
    \item 采用2位并行乘法方式，实现硬件开销小、实现最简单的优化方案。
    \item 完成功能仿真与上板验证，确保优化后乘法器结果正确。
    \item 对比优化前后性能，完成实验分析与总结。
\end{enumerate}

\section{实验原理图}
\subsection{2位并行乘法器结构原理图}
本实验采用2位并行乘法优化方案，每次同时读取乘数低2位，根据不同组合生成0倍、1倍、2倍、3倍被乘数，作为部分积参与累加，从而将时钟周期减少一半。

\begin{figure}[H]
    \centering
    \includegraphics[width=0.8\textwidth]{原理图.jpg}
    \caption{2位并行乘法器硬件结构图}
    \label{fig:mult}
\end{figure}

\subsection{原理说明}
\begin{enumerate}
    \item \textbf{符号与绝对值处理单元}：提取操作数符号位，对负数进行取反加1得到绝对值。
    \item \textbf{被乘数寄存器}：64位，每次左移2位。
    \item \textbf{乘数寄存器}：32位，每次右移2位。
    \item \textbf{2位编码与部分积生成单元}：根据乘数低两位选择部分积（0、×1、×2、×3）。
    \item \textbf{累加寄存器}：对部分积进行累加。
    \item \textbf{控制单元}：判断乘数是否为0，产生乘法结束信号。
    \item \textbf{符号修正单元}：根据符号位异或结果，对最终乘积进行符号修正。
\end{enumerate}

\section{实验过程与代码实现}
本次实验仅修改 \texttt{multiply.v} 文件，保持其他模块不变。

\subsection{核心优化思路}
\begin{itemize}
    \item 原始乘法：每次处理1位，需要32个时钟周期。
    \item 优化乘法：每次处理2位，只需要16个时钟周期。
    \item 部分积直接由移位和加法生成，硬件最简单、最容易实现。
\end{itemize}

\subsection{优化后Verilog代码}
\begin{lstlisting}
`timescale 1ns / 1ps

module multiply(              // 乘法器
    input         clk,        // 时钟
    input         mult_begin, // 乘法开始信号
    input  [31:0] mult_op1,   // 乘法源操作数1
    input  [31:0] mult_op2,   // 乘法源操作数2
    output [63:0] product,    // 乘积
    output        mult_end    // 乘法结束信号
);

    //乘法正在运算信号和结束信号
    reg mult_valid;
    assign mult_end = mult_valid & ~(|multiplier); //乘法结束信号：乘数全0
    always @(posedge clk)
    begin
        if (!mult_begin || mult_end)
        begin
            mult_valid <= 1'b0;
        end
        else
        begin
            mult_valid <= 1'b1;
        end
    end

    //两个源操作取绝对值
    wire        op1_sign;      //操作数1的符号位
    wire        op2_sign;      //操作数2的符号位
    wire [31:0] op1_absolute;  //操作数1的绝对值
    wire [31:0] op2_absolute;  //操作数2的绝对值
    assign op1_sign = mult_op1[31];
    assign op2_sign = mult_op2[31];
    assign op1_absolute = op1_sign ? (~mult_op1+1) : mult_op1;
    assign op2_absolute = op2_sign ? (~mult_op2+1) : mult_op2;

    //被乘数，每次左移2位
    reg  [63:0] multiplicand;
    always @ (posedge clk)
    begin
        if (mult_valid)
        begin    
            multiplicand <= {multiplicand[61:0],2'b0};
        end
        else if (mult_begin) 
        begin   
            multiplicand <= {32'd0,op1_absolute};
        end
    end

    //乘数，每次右移2位
    reg  [31:0] multiplier;
    always @ (posedge clk)
    begin
        if (mult_valid)
        begin   
            multiplier <= {2'b0,multiplier[31:2]}; 
        end
        else if (mult_begin)
        begin   
            multiplier <= op2_absolute; 
        end
    end

    // 2位并行部分积生成
    wire [63:0] partial_product;
    assign partial_product = (multiplier[1:0] == 2'b00) ? 64'd0 : 
                            (multiplier[1:0] == 2'b01) ? multiplicand : 
                            (multiplier[1:0] == 2'b10) ? (multiplicand << 1) : 
                            (multiplicand + (multiplicand << 1));

    //累加器
    reg [63:0] product_temp;
    always @ (posedge clk)
    begin
        if (mult_valid)
        begin
            product_temp <= product_temp + partial_product;
        end
        else if (mult_begin) 
        begin
            product_temp <= 64'd0;  
        end
    end 

    //结果符号处理
    reg product_sign;
    always @ (posedge clk)
    begin
        if (mult_valid)
        begin
              product_sign <= op1_sign ^ op2_sign;
        end
    end 
    assign product = product_sign ? (~product_temp+1) : product_temp;
endmodule
\end{lstlisting}

\section{实验结果与分析}
\subsection{仿真波形图}
\begin{figure}[H]
    \centering
    \includegraphics[width=0.9\textwidth]{波形图.jpg}
    \caption{乘法器仿真波形图}
\end{figure}

\textbf{波形说明：}
\begin{itemize}
    \item 乘法启动信号 \texttt{mult\_begin} 有效后，乘法器开始工作。
    \item 乘数每次右移2位，被乘数每次左移2位，周期数大幅减少。
    \item 16个时钟周期后 \texttt{mult\_end} 拉高，表示乘法完成。
    \item 输出乘积 \texttt{product} 与理论计算值完全一致，结果正确。
\end{itemize}

\subsection{优化前后时钟周期对比}
\begin{table}[H]
    \centering
    \begin{tabular}{|c|c|}
        \hline
        乘法器类型 & 完成32位乘法所需时钟周期 \\
        \hline
        原始串行乘法器 & 32 周期 \\
        \hline
        2位并行优化乘法器 & 16 周期 \\
        \hline
    \end{tabular}
    \caption{性能对比表}
\end{table}

\textbf{结论：}
优化后乘法器时钟周期减少50\%，性能提升一倍，硬件实现简单，资源开销小。

\subsection{3. 上板验证结果}
\begin{figure}[H]
    \centering
    \includegraphics[width=0.8\textwidth]{1.jpg}
    \caption{实验箱上板验证结果1}
\end{figure}
\begin{figure}[H]
    \centering
    \includegraphics[width=0.8\textwidth]{2.jpg}
    \caption{实验箱上板验证结果2}
\end{figure}

\textbf{结果说明：} \\
多次测试均正确，优化后乘法器功能完全正常。

\section{实验收获与感想}
通过本次定点乘法器优化实验，我收获如下：
\begin{enumerate}
    \item 理解了原始串行乘法器需要逐位计算，周期长、效率低。
    \item 掌握了最简单有效的2位并行乘法优化方法，只需修改移位位数和部分积判断，即可将周期减半。
    \item 学会了部分积的生成方式：使用移位和加法实现×1、×2、×3，无需乘法器，硬件更简洁高效。
    \item 完成了从代码编写、功能仿真到上板验证的完整流程，提升了硬件设计能力。
    \item 认识到计算机组成原理中“并行化”是提升运算速度的核心思想，简单优化即可带来明显性能提升。
\end{enumerate}

\section{实验总结}
本次实验成功完成了32位定点乘法器的性能优化，采用2位并行乘法方式，将运算周期从32个时钟周期减少至16个时钟周期，性能提升一倍。
优化后的乘法器逻辑简单、易于实现、资源占用少，仿真与上板验证均证明结果正确。

\end{document}